\section{Vector boundary information for a moving-support local family}
\label{sec:v22-vector-information}

We formulate the moving-support experiment without assuming absolute
continuity with respect to its reference member.  The original family is
jointly dominated by Lebesgue measure and its failure atom, but a likelihood
ratio relative to the reference law need not exist on the support-exclusive
part.  We separate that mass and apply parameter-independent common-collar
censoring.  The censored family is dominated by its censored reference member
and satisfies a uniform LAN expansion.  The original family is uniformly
asymptotically equivalent to that representative on every compact local set.
We distinguish this equivalence statement from finite-subexperiment
convergence to the Gaussian shift and formulate singular information on its
identifiable quotient.  Proposition~\ref{prop:v26-common-domination} records
the common measures explicitly.

Throughout this section we use the Hellinger convention
\begin{equation}\label{eq:v22-hellinger-convention}
 H^2(P,Q)=\int(\sqrt p-\sqrt q)^2\,d\mu
        =2-2\int\sqrt{pq}\,d\mu,
\end{equation}
for any common dominating measure $\mu$.  With this convention
$H^2(P^{\otimes k},Q^{\otimes k})\le kH^2(P,Q)$.

Let $\Theta\subset\mathbb R^r$ be a neighborhood of the origin.  On a fixed
bounded open set $U\subset\mathbb R^m$ suppose
\begin{equation}\label{eq:v22-vector-density}
 f_\theta(x)=a_\theta(x)(w_\theta(x))_+,
 \qquad \int_U f_\theta(x)\,dx=1.
\end{equation}
Assume that $a_\theta,w_\theta$ are $C^3$ jointly in $(\theta,x)$ on a
common compact neighborhood of the relevant supports, uniformly for small
$\theta$, that $a_\theta$ is bounded above and bounded away from zero there,
and that
\[
 \Sigma=\{w_0=0\}\cap\operatorname{supp}a_0
\]
is a compact regular hypersurface with $|\nabla w_0|\ge c_*>0$.  After
shrinking $\Theta$, all moving boundaries lie in one fixed tubular collar.
Put
\begin{equation}\label{eq:v22-boundary-velocity-matrix}
 V(x)=D_\theta w_\theta(x)|_{\theta=0}\in\mathbb R^r,
 \qquad
 \mathcal J_\Sigma=
 \int_\Sigma\frac{a_0(x)V(x)V(x)^t}{|\nabla w_0(x)|}\,d\sigma(x).
\end{equation}
Multiplying $w_\theta$ by a positive smooth factor and dividing
$a_\theta$ by the same factor leaves $\mathcal J_\Sigma$ unchanged.

Let $\delta_n\downarrow0$, $0<p_n\le1$, and assume
\begin{equation}\label{eq:v22-vector-rate}
 n p_n\delta_n^2\log(1/\delta_n)\longrightarrow1.
\end{equation}
For $h\in\mathbb R^r$ define on
$\Omega=U\sqcup\{\dagger\}$
\begin{equation}\label{eq:v22-thinned-family}
 P_{n,h}=(1-p_n)\delta_\dagger+p_nf_{\delta_nh}(x)\,dx.
\end{equation}
No absolute continuity of $P_{n,h}$ with respect to $P_{n,0}$ is assumed.

\subsection{Collar estimates and the singular part}

\begin{lemma}[Uniform vector collar estimates]
\label{lem:v22-vector-collar}
Let
\[
 \mathsf S(x)=D_\theta\log a_\theta(x)|_0+\frac{V(x)}{w_0(x)},
 \qquad D=\{w_0>0\}.
\]
There are $q_0,C>0$ such that, for $0<q<q_0$ and unit vectors
$u,v\in\mathbb R^r$,
\begin{align}
 \mathbb E_0[(u^t\mathsf S)(v^t\mathsf S)
       \mathbf1_{\{w_0\ge q\}}]
 &=u^t\mathcal J_\Sigma v\log(1/q)+O(1),
                                      \label{eq:v22-vector-second}\\
 \mathbb E_0[|u^t\mathsf S|^3\mathbf1_{\{w_0\ge q\}}]
 &\le Cq^{-1},                       \label{eq:v22-vector-third}\\
 \mathbb E_0[|u^t\mathsf S|^4\mathbf1_{\{w_0\ge q\}}]
 &\le Cq^{-2},                       \label{eq:v22-vector-fourth}\\
 \left\|\mathbb E_0[\mathsf S\mathbf1_{\{w_0\ge q\}}]\right\|
 &\le Cq,
 \qquad
 \mathbb P_0(0<w_0<q)\le Cq^2.      \label{eq:v22-vector-centering}
\end{align}
The constants are uniform for the indicated unit directions and on compact
families satisfying the displayed regularity bounds.
\end{lemma}

\begin{proof}
Use normal coordinates $(y,s)$ with $s=w_0$ in a fixed tubular neighborhood
of $\Sigma$.  The Jacobian is $|\nabla w_0|^{-1}+O(s)$ and the null density
is $a_0s$.  The only nonintegrable contribution to the second score moment is
$a_0(u^tV)(v^tV)/s$, whose normal integral gives
\eqref{eq:v22-vector-second}.  The $k$th singular moment has normal integrand
$O(s^{1-k})$, proving \eqref{eq:v22-vector-third} and
\eqref{eq:v22-vector-fourth}.  The same coordinates give the quadratic collar
mass.  Differentiating $\int f_\theta=1$ at zero gives
$\mathbb E_0\mathsf S=0$ in the improper integrable sense; the omitted collar
has first score moment $O(q)$.
\end{proof}

\begin{lemma}[Support-exclusive mass]
\label{lem:v22-support-decomposition}
For every compact $K\subset\mathbb R^r$ there is $C_K$ such that, uniformly
for $h\in K$ and all large $n$,
\begin{align}
 P_{n,h}^{\perp0}(\Omega)&\le C_Kp_n\delta_n^2,
                                      \label{eq:v22-alt-singular-mass}\\
 P_{n,0}\{f_{\delta_nh}=0<f_0\}&\le C_Kp_n\delta_n^2.
                                      \label{eq:v22-null-exclusive-mass}
\end{align}
Consequently, under either product experiment, the probability of seeing at
least one support-exclusive observation is at most
$C_Knp_n\delta_n^2=o(1)$.
\end{lemma}

\begin{proof}
Joint $C^1$ regularity gives
$|w_{\delta_nh}-w_0|\le C_K\delta_n$ on the common collar.  A point belonging
to exactly one support therefore lies in $|w_0|\le C_K\delta_n$.  On that
strip the relevant density is $O(\delta_n)$ and the strip has volume
$O(\delta_n)$.  Multiplication by $p_n$ gives the one-observation estimates.
A union bound and \eqref{eq:v22-vector-rate} give the product estimate.
\end{proof}

\subsection{A reference-dominated common-collar representative}

Put
\begin{equation}\label{eq:v22-truncation-scale}
 L_n=(\log(1/\delta_n))^{1/4},\qquad
 q_n=\delta_nL_n,\qquad C_n=\{w_0\ge q_n\}.
\end{equation}
Introduce a second cemetery symbol $\partial$ and let
$T_n:\Omega\to C_n\sqcup\{\dagger,\partial\}$ keep $\dagger$ and every
$x\in C_n$, while sending every successful $x\notin C_n$ to $\partial$.
Write
\begin{equation}\label{eq:v22-censored-law}
 \overline P_{n,h}=T_nP_{n,h}
 =(1-p_n)\delta_\dagger
 +p_nf_{\delta_nh}\mathbf1_{C_n}\,dx
 +p_nr_{n,h}\delta_\partial,
 \qquad
 r_{n,h}=\int_{C_n^c}f_{\delta_nh}\,dx.
\end{equation}

\begin{lemma}[Common-collar asymptotic equivalence]
\label{lem:v22-common-collar-equivalence}
For every compact $K\subset\mathbb R^r$,
\begin{equation}\label{eq:v22-collar-mass}
 \sup_{h\in K}r_{n,h}\le C_Kq_n^2,
 \qquad
 \sup_{h\in K}np_nr_{n,h}=o(1).
\end{equation}
For all large $n$, $\{\overline P_{n,h}:h\in K\}$ is dominated by
$\overline P_{n,0}$, and
\begin{equation}\label{eq:v22-lecam-censoring}
 \Delta\!\left(
  \{P_{n,h}^{\otimes n}:h\in K\},
  \{\overline P_{n,h}^{\otimes n}:h\in K\}
 \right)
 \le C_Knp_nq_n^2
 =O((\log(1/\delta_n))^{-1/2})=o(1).
\end{equation}
Both comparison kernels can be chosen independently of $h$.
\end{lemma}

\begin{proof}
For $h\in K$, the moving support lies in
$\{w_0>-C_K\delta_n\}$.  Integrating
$f_{\delta_nh}=O(w_0+C_K\delta_n)$ over
$-C_K\delta_n<w_0<q_n$ gives $r_{n,h}=O(q_n^2)$.  Equation
\eqref{eq:v22-vector-rate} gives the second bound.

On $C_n$,
$w_{\delta_nh}\ge q_n-C_K\delta_n\ge q_n/2$ for all large $n$, uniformly on
$K$, so all successful densities are mutually positive there.  Moreover
$r_{n,0}>0$: the null collar $0<w_0<q_n$ has positive mass.  Hence the
$\partial$ atom of the null censored law has positive mass, and the whole
censored family is dominated by $\overline P_{n,0}$.

The forward kernel is the parameter-independent censoring map $T_n$.  For a
reverse kernel, keep $C_n$ and $\dagger$ fixed and send $\partial$ to one
fixed point of the original successful sample space.  This reverse kernel is
also parameter independent.  Couple it with the original observation before
censoring; disagreement is possible only if censoring occurred.  A union
bound gives probability at most $np_nr_{n,h}$, uniformly on $K$, proving
\eqref{eq:v22-lecam-censoring}.
\end{proof}

\subsection{Uniform LAN on the representative}

\begin{lemma}[Uniform LAN]
\label{lem:v22-dominated-lan}
For $h$ in a fixed compact set define under
$\overline P_{n,0}^{\otimes n}$
\begin{align}
 \Delta_n&=\delta_n\sum_{i=1}^nY_i
 \mathsf S(X_i)\mathbf1_{C_n}(X_i),
                                      \label{eq:v22-delta}\\
 Q_n&=\delta_n^2\sum_{i=1}^nY_i
 \mathsf S(X_i)\mathsf S(X_i)^t\mathbf1_{C_n}(X_i),
                                      \label{eq:v22-Q}
\end{align}
where $Y_i$ is the success indicator and, conditionally on success under the
null, $X_i$ has density $f_0$.  Then
\begin{equation}\label{eq:v22-dominated-expansion}
 \log\frac{d\overline P_{n,h}^{\otimes n}}
          {d\overline P_{n,0}^{\otimes n}}
 =h^t\Delta_n-\frac12h^tQ_nh+o_{\overline P_{n,0}}(1)
\end{equation}
uniformly for $h$ in compact subsets of $\mathbb R^r$, and
\begin{equation}\label{eq:v22-central-limits}
 \Delta_n\Rightarrow N_r(0,\mathcal J_\Sigma),
 \qquad Q_n\longrightarrow\mathcal J_\Sigma
 \quad\text{in probability}.
\end{equation}
The probability that the atom $\partial$ occurs in the product experiment is
$o(1)$ uniformly on compact local parameter sets.
\end{lemma}

\begin{proof}
The last assertion is Lemma~\ref{lem:v22-common-collar-equivalence}.  On its
complement the likelihood ratio uses only the common support $C_n$ and the
unchanged failure atom.  Taylor expansion before division by $w_0$ gives,
uniformly for bounded $h$ and $x\in C_n$,
\[
 \log\frac{f_{\delta_nh}(x)}{f_0(x)}
 =\delta_nh^t\mathsf S(x)
 -\frac{\delta_n^2}{2}(h^t\mathsf S(x))^2+R_{n,h}(x).
\]
The regular parameter remainder has expectation $O(\delta_n^2)$ per
successful observation, while the cubic logarithmic remainder is bounded in
expectation by
$C_K\delta_n^3\mathbb E_0(1+\|\mathsf S\|^3)\mathbf1_{C_n}$.
Lemma~\ref{lem:v22-vector-collar} and \eqref{eq:v22-vector-rate} give
\[
 np_n\delta_n^2=o(1),\qquad
 np_n\delta_n^3q_n^{-1}=o(1),
\]
which proves \eqref{eq:v22-dominated-expansion}.

The truncated-score mean has norm at most $Cq_n$, so
$\|\mathbb E\Delta_n\|\le Cnp_n\delta_nq_n=o(1)$.  Its covariance tends to
$\mathcal J_\Sigma$ because
$\log(1/q_n)/\log(1/\delta_n)\to1$.  On $C_n$ every scalar score contribution
multiplied by $\delta_n$ is $O(L_n^{-1})$; Lindeberg therefore holds uniformly
on the unit sphere and Cram\'er--Wold gives the first limit in
\eqref{eq:v22-central-limits}.  Finally
\[
 np_n\delta_n^4
 \mathbb E_0[\|\mathsf S\|^4\mathbf1_{C_n}]
 \le Cnp_n\delta_n^4q_n^{-2}=o(1),
\]
while the mean of $Q_n$ tends to $\mathcal J_\Sigma$.  Hence
$Q_n\to\mathcal J_\Sigma$ in probability.
\end{proof}

\begin{lemma}[Bounded-sequence contiguity]
\label{lem:v22-contiguity}
If $h_n$ is any bounded sequence, then
$\overline P_{n,h_n}^{\otimes n}$ and
$\overline P_{n,0}^{\otimes n}$ are mutually contiguous.  The same is true
for $P_{n,h_n}^{\otimes n}$ and $P_{n,0}^{\otimes n}$.
\end{lemma}

\begin{proof}
Along every subsequence, take a further subsequence on which $h_n\to h$.
Lemma~\ref{lem:v22-dominated-lan} gives joint convergence of the log
likelihood ratio under the null to
$h^tZ-\frac12h^t\mathcal J_\Sigma h$ with
$Z\sim N(0,\mathcal J_\Sigma)$.  Its exponential has mean one, so Le Cam's
first lemma gives contiguity in one direction; the same argument under the
local alternative, or the third lemma applied to $-h$, gives the reverse
direction.  Since every subsequence has this property, the whole bounded
sequence is mutually contiguous.  Lemma
\ref{lem:v22-common-collar-equivalence} transfers contiguity to the original
experiments.
\end{proof}

\subsection{The Gaussian limit on the identifiable quotient}

Let
\begin{equation}\label{eq:v22-identifiable-space}
 \mathsf H=\operatorname{Ran}(\mathcal J_\Sigma),\qquad
 \Pi_{\mathsf H}:\mathbb R^r\to\mathsf H
\end{equation}
be the orthogonal projection.  The restriction
$\mathcal J_{\mathsf H}=\mathcal J_\Sigma|_{\mathsf H}$ is positive definite.
The Gaussian limit experiment is most cleanly written in central-sequence
coordinates as
\begin{equation}\label{eq:v22-gaussian-central-experiment}
          \Delta\sim N_{\mathsf H}(\mathcal J_{\mathsf H}
                   \Pi_{\mathsf H}h,\mathcal J_{\mathsf H}).
\end{equation}
Equivalently,
$Z=\Pi_{\mathsf H}h+\mathcal J_{\mathsf H}^{-1/2}G$ with standard Gaussian
$G$ on $\mathsf H$.  Directions in $\ker\mathcal J_\Sigma$ are not
identifiable at the critical logarithmic scale.

\begin{theorem}[Moving-support vector boundary Gaussian limit]
\label{thm:v22-vector-boundary-gaussian}
Under \eqref{eq:v22-vector-density}--\eqref{eq:v22-vector-rate}:

\begin{enumerate}
\item On every compact local parameter set, the original experiment is at
$o(1)$ Le Cam distance from the reference-dominated common-collar representative, by
\eqref{eq:v22-lecam-censoring}.  The representative satisfies the compact-
uniform LAN expansion \eqref{eq:v22-dominated-expansion} and the central
limits \eqref{eq:v22-central-limits}.

\item For every finite collection $h_1,\ldots,h_m$, the corresponding finite
local subexperiment converges in Le Cam distance to the restriction of the
Gaussian experiment \eqref{eq:v22-gaussian-central-experiment} to those
parameters.  More generally, every bounded local sequence satisfies the
contiguity and third-lemma conclusions implied by the LAN expansion.  We do
not identify this finite-subexperiment statement with a separate assertion
of global deficiency convergence for an uncountable parameter set.

\item For every fixed $h$,
\begin{equation}\label{eq:v22-testing-profile}
 \|P_{n,h}^{\otimes n}-P_{n,0}^{\otimes n}\|_{\rm TV}
 \longrightarrow
 2\Phi\!\left(\frac{\sqrt{h^t\mathcal J_\Sigma h}}2\right)-1.
\end{equation}
If $\mathcal J_\Sigma$ is positive definite, the reference-model statistic
\begin{equation}\label{eq:v22-local-estimator}
                 \widehat h_n=\mathcal J_\Sigma^{-1}\Delta_n
\end{equation}
(with observations outside $C_n$ contributing zero) satisfies
$\widehat h_n\Rightarrow N_r(h,\mathcal J_\Sigma^{-1})$ uniformly on compact
local sets.  If $\mathcal J_\Sigma$ is singular, the corresponding estimator
of the identifiable coordinate is
$\widehat h_n^{\rm id}=\mathcal J_\Sigma^+\Delta_n$, with limit centered at
$\Pi_{\mathsf H}h$ and covariance $\mathcal J_\Sigma^+$.

\item Let $W$ be positive semidefinite on $\mathsf H$ and
$\ell_M(z)=\min\{z^tWz,M\}$.  The standard local asymptotic minimax theorem
applied to the LAN experiment on $\mathsf H$ gives
\begin{equation}\label{eq:v22-minimax-truncated}
 \lim_{R\to\infty}\liminf_{n\to\infty}
 \inf_{T_n}\sup_{\substack{h\in\mathsf H\\\|h\|\le R}}
 \mathbb E_{n,h}\ell_M(T_n-h)
 \ge
 \mathbb E\ell_M(Z_0),
 \qquad
 Z_0\sim N_{\mathsf H}(0,\mathcal J_{\mathsf H}^{-1}).
\end{equation}
Letting $M\to\infty$ gives
$\operatorname{tr}(W\mathcal J_{\mathsf H}^{-1})$.  When
$\mathcal J_\Sigma$ is positive definite, this is the full-parameter bound.
\end{enumerate}
\end{theorem}

\begin{proof}
Item 1 is Lemmas~\ref{lem:v22-common-collar-equivalence} and
\ref{lem:v22-dominated-lan}.  For any finite parameter set, the joint
likelihood-ratio vector is a continuous function of $\Delta_n,Q_n$ plus a
uniform $o_P(1)$ term.  Equations \eqref{eq:v22-central-limits} therefore give
convergence of the likelihood-ratio experiment to
\eqref{eq:v22-gaussian-central-experiment}; the finite-parameter version of
Le Cam's convergence theorem gives the asserted Le Cam-distance convergence.
Lemma~\ref{lem:v22-contiguity} gives the bounded-sequence statement.

Le Cam's third lemma translates the limiting central sequence by
$\mathcal J_\Sigma h$.  The equal-covariance Gaussian total-variation formula
gives \eqref{eq:v22-testing-profile}.  The estimator statements follow by
linear transformation of the central sequence; in the singular case the
Moore--Penrose inverse acts only on $\mathsf H$.

For compact local alternatives the likelihood ratio on $C_n$ is bounded above
and below by constants tending to one, since $\delta_n/w_0\le L_n^{-1}$.
The third- and fourth-moment bounds therefore remain uniform under local
alternatives.  The standard fourth-moment formula for sums and
\eqref{eq:v22-vector-second}--\eqref{eq:v22-vector-fourth} give
$\sup_{h\in K}\mathbb E_{n,h}\|\Delta_n\|^4\le C_K$; the censored event is
$o(1)$ and contributes zero to the displayed statistics.  This gives uniform
integrability for compact-local quadratic risk.  Finally apply the classical
finite-dimensional LAN local asymptotic minimax theorem on the identifiable
space $\mathsf H$; see \cite{IbragimovHasminskii,HiranoPorter}.
\end{proof}

\begin{remark}[Reference-model nature of the central sequence]
The statistics \eqref{eq:v22-delta} and \eqref{eq:v22-local-estimator} use the
score and information at the reference geometry.  They are local-experiment
statistics, not globally adaptive estimators for an unknown reference model.
The physical theorem below uses them only after a reference channel and a
finite local chart have been fixed.
\end{remark}

\subsection{Quantitative stability of the moving boundary}

\begin{lemma}[Linearly vanishing support stability]
\label{lem:v22-hypersurface-stability}
Let $f=A(w)_+$ and $\widetilde f=\widetilde A(\widetilde w)_+$ be normalized
densities on the same bounded open set.  Suppose $A,\widetilde A$ are
uniformly positive and $C^1$, the zero sets of $w,\widetilde w$ are uniformly
regular in one fixed tubular collar, and
\[
 \|w-\widetilde w\|_{C^1}\le\varepsilon<1/2,
 \qquad
 \|A-\widetilde A\|_{C^0}\le\eta.
\]
Then, with the convention \eqref{eq:v22-hellinger-convention},
\begin{equation}\label{eq:v22-hypersurface-stability}
 H^2(f\,dx,\widetilde f\,dx)
 \le C\left\{\varepsilon^2\log\frac e\varepsilon+\eta^2\right\}.
\end{equation}
The estimate is unchanged, up to the uniform constant, if a common positive
smooth factor is moved between the amplitude and defining function.
\end{lemma}

\begin{proof}
Use the tubular coordinate $s=w$.  On $|s|\le C\varepsilon$ the total mass of
both densities is $O(\varepsilon^2)$.  On $s>C\varepsilon$ the supports agree
and
\[
 |\sqrt{Aw}-\sqrt{\widetilde A\widetilde w}|^2
 \le C\left(\eta^2s+\frac{\varepsilon^2}{s}\right).
\]
The first term is $O(\eta^2)$ after integration and the second is
$O(\varepsilon^2\log(e/\varepsilon))$.  The same estimate holds using
$\widetilde w$ as the normal coordinate on the other side of the collar;
away from the collar all denominators are uniformly bounded.  Multiplication
of $w,\widetilde w$ by the same uniformly positive smooth factor merely
changes the tubular Jacobian and amplitude bounds by uniform constants.
\end{proof}

\begin{remark}[Terminology]
We reserve ``LAN expansion'' for the reference-dominated common-collar representative.
For the original family we state asymptotic equivalence to that representative
and the resulting finite-subexperiment Gaussian limit, testing profile,
contiguity, and minimax consequences.  The original family already has a common dominating measure.
Reference absolute continuity, local asymptotic normality, and convergence
of experiments are distinct assertions.
\end{remark}
